\documentclass[11pt]{article}
\usepackage[T1]{fontenc}
\usepackage{lmodern,amsmath,amssymb,amsthm,mathtools,microtype,xcolor,enumitem}
\usepackage[letterpaper,left=1.02in,right=1.02in,top=0.94in,bottom=0.94in]{geometry}
\usepackage[colorlinks=true,linkcolor=blue!45!black,urlcolor=blue!45!black,citecolor=blue!45!black]{hyperref}
\setlist{nosep}
\setlength{\emergencystretch}{1.5em}
\newtheorem{theorem}{Theorem}[section]
\newtheorem{proposition}[theorem]{Proposition}
\newtheorem{lemma}[theorem]{Lemma}
\newtheorem{corollary}[theorem]{Corollary}
\newcommand{\R}{\mathbb R}
\newcommand{\C}{\mathbb C}
\newcommand{\dd}{\,\mathrm d}
\newcommand{\ol}{\overline}
\DeclareMathOperator{\diag}{diag}
\title{Zero multiplicities from three nearby prime sums}
\author{Research note}
\date{27 November 2026}
\begin{document}
\maketitle
\begin{abstract}
Under RH, simple zeros of zeta have density one if and only if the mean square of a second difference of a smoothed prime sum, at dyadic lags $h$, is asymptotic to $(2/3)h^3\log(1/h)$. At arbitrary finite limiting coefficients we prove a converse for every three-point geometry. The same observations measure multiplicity-weighted common zeros of Dirichlet $L$-functions. Attaining the minimum coefficient in the single class $1\bmod q$ forces the entire limiting residue-class covariance matrix to be scalar. For the full matrix on several nearby scales we prove the order of each eigenvalue, allowing arbitrary coincidences of zeros. We also reduce the factor $4/3$ in the source article's zeta coalescence bounds to $1.3208$.
\end{abstract}

\section{Reading multiplicity from a prime sum}
Let $\Lambda(p^k)=\log p$ for primes $p$ and $k\ge1$, with $\Lambda(n)=0$ otherwise. For $t\ge0$ write
\[
 E(t)=e^{t/2}\left[\sum_{n\le e^t}\Lambda(n)(1/n-e^{-t})-t+1+\gamma_{\rm E}\right],
\]
where $\gamma_{\rm E}$ is Euler's constant. Fix $a,b>0$. The contrast
\[
 D_h f(t)=b f(t)-(a+b)f(t+ah)+a f(t+(a+b)h)
\]
annihilates affine functions of $t$. Put $\tau(h)=h^3\log(1/h)$ for $0<h<1$ and $B_{a,b}=a^2b^2(a+b)/3$. Define
\[
 V_{a,b}(h)=\lim_{Y\to\infty}Y^{-1}\int_0^Y |D_hE(t)|^2\dd t,
 \qquad M_2(T)=\sum_{0<\gamma\le T}^{\rm distinct}m(\gamma)^2 .
\]
Here $m(\gamma)$ is the multiplicity at $1/2+i\gamma$. Write $N(T)$ for the positive-ordinate count up to $T$ with multiplicity.

\begin{theorem}\label{th:zeta}
Assume RH. For every fixed $a,b>0$ and every finite $\kappa\ge0$,
\[
 \lim_{h\downarrow0}\frac{V_{a,b}(h)}{\tau(h)}=B_{a,b}\kappa
 \quad\Longleftrightarrow\quad
 \lim_{T\to\infty}\frac{2\pi M_2(T)}{T\log T}=\kappa .
\]
It suffices on the left to use $h=1/n$ for integers $n\to\infty$. At $\kappa=1$ it suffices to use $h=2^{-n}$, and both conditions are equivalent to density-one simplicity of zeta zeros.
\end{theorem}

For equally spaced nodes $a=b=1$, the measured coefficient is $2\kappa/3$. Knowing that one such limit exists gives the limiting coefficient for every three-point geometry. The article \cite{repo} establishes the forward variance formula and asks whether its factor $4/3$ for the upper bound is optimal, and what happens for the full residue-class matrix. We give a converse, a smaller upper bound, and the orders of every eigenvalue of the full residue-class covariance matrix. We also characterize all finite limits for the same contrast in each coordinate. Our density-one criterion does not determine whether it is satisfied. The average in $t$ is always taken first.

\section{An inverse theorem for positive measures}
Write
\[
 P(x)=b-(a+b)e^{iax}+ae^{i(a+b)x},\qquad K(x)=|P(x)|^2/x^4
\]
with $K(0)$ given by continuity. The cancellations give $|P(x)|\ll_{a,b}\min(x^2,1)$ and $K(0)=a^2b^2(a+b)^2/4>0$. Let $A$ be a positive locally finite measure on $[0,\infty)$ with $\int_1^\infty x^{-4}\dd A(x)<\infty$. Put $A(X)=A([0,X])$ and $U(h)=h^4\int K(hx)\dd A(x)$.

\begin{theorem}\label{th:inverse}
For finite $d\ge0$,
\[
 \frac{\pi A(T)}{T\log T}\longrightarrow d
 \quad\Longleftrightarrow\quad
 \frac{U(h)}{\tau(h)}\longrightarrow d B_{a,b}.
\]
For the reverse implication it suffices that the asserted limit hold on any decreasing sequence $h_n\to0$ with $h_{n+1}/h_n\to1$.
\end{theorem}

\begin{proof}
Set $\mu_T(E)=\pi A(TE)/(T\log T)$ for $T>1$. If the counting limit holds, $A(X)\ll(1+X)\log(2+X)$ and $\mu_T$ converges vaguely on $[0,\infty)$ to $d\,\mathrm dx$ by convergence of the distribution functions at every positive $x$ and vanishing mass near zero. For $s>0$,
\begin{equation}\label{eq:scale}
 \frac{U(s/T)}{T^{-3}\log T}
  =\frac1\pi\int_0^\infty |P(sx)|^2x^{-4}\dd\mu_T(x).
\end{equation}
For $T\ge2$, $R\ge1$ the growth bound gives $\mu_T([0,R])\ll R(1+\log(2R))$. A dyadic sum bounds the integral over $x>R$ by $O(R^{-3}(1+\log(2R)))$. The integrands are continuous at zero and locally uniformly continuous in $(s,x)$. Consequently vague convergence suffices in \eqref{eq:scale}, even if $s$ tends to a fixed positive number.

For the integral constant let $u=(0,a,a+b)$, $c=(b,-(a+b),a)$, and $H(t)=\sum_i c_i(t-u_i)_+$. This tent has base $a+b$, height $ab$, and vanishes off $[0,a+b]$. From $H''=\sum c_i\delta_{u_i}$, Plancherel in the angular-frequency convention gives
\begin{equation}\label{eq:tent}
 \frac1\pi\int_0^\infty K(x)\dd x=\int_{\R}H(t)^2\dd t=B_{a,b}.
\end{equation}
Taking $s=1$ proves the forward direction.

Conversely $K$ is bounded below by a positive constant on $[0,\delta]$ for some $\delta>0$. A bound $U(h_n)\ll\tau(h_n)$ gives $A(\delta/h_n)\ll h_n^{-1}\log(1/h_n)$. Monotonicity and the ratio assumption give $A(X)\ll (1+X)\log(2+X)$ at all $X$. Thus the measures $\mu_T$ have vaguely convergent subsequences along every sequence $T\to\infty$. Every limit $\nu$ satisfies $\nu([0,R])\le C R$ for all $R>0$ with the same $C$. Indeed, apply the growth estimate at $(R+\epsilon)T$ to a continuous majorant of the indicator of $[0,R]$ supported in $[0,R+\epsilon]$, bounded by one, and send $\epsilon$ to zero. In particular $\nu(\{0\})=0$.

Along such a subsequence, for any $s>0$ choose $n=n(T)$ with $T h_n\to s$. The tail and continuity bounds above give
\[
 \int |P(sx)|^2x^{-4}\dd\nu(x)=\pi d B_{a,b}s^3 .
\]
By \eqref{eq:tent} the measure $d\,\mathrm dx$ gives the same value. Write $\sigma=\nu-d\,\mathrm dx$, whose total variation on $[0,R]$ is $O(R)$ for all $R>0$. Define
\[
 R_\sigma(t)=\int_0^\infty (\cos(tx)-1+(tx)^2/2)x^{-4}\dd\sigma(x).
\]
The bound $|\cos y-1+y^2/2|\ll\min(y^4,y^2)$, splitting at $1/|t|$ and summing the tail dyadically give $|R_\sigma(t)|\ll|t|^3$ on the whole line. Expanding $|P|^2$ cancels the constant and quadratic terms and gives
\[
 ab R_\sigma((a+b)s)=b(a+b)R_\sigma(as)+a(a+b)R_\sigma(bs).
\]
Let $M=\sup_{t>0}|R_\sigma(t)|/t^3<\infty$. Dividing by $ab(a+b)^3s^3$ gives
\[
 M\le \frac{a^2+b^2}{(a+b)^2}\,M.
\]
The coefficient is strictly less than one, so $R_\sigma$ vanishes.

Its fourth distributional derivative is the cosine transform of $\sigma$. To justify differentiation, truncate the measure to compact subintervals of $(0,\infty)$ and integrate by parts four times against a Schwartz test function. The common bound $O(|t|^3)$ gives convergence on the original side, and rapid decay of the transformed test function gives convergence on the derivative side. The even extension of $\sigma$ has zero Fourier transform and hence vanishes. All subsequential limits therefore equal $d\,\mathrm dx$. Applying convergence to $[0,1]$ proves the counting limit.
\end{proof}

\paragraph{Changing the denominator.}
In the applications $x^{-4}$ occurs as $(1/4+x^2)^{-2}$. For a measure with $\int(1+x^2)^{-1}\dd A(x)<\infty$ the change costs $O_{A,a,b}(h^4)=o(\tau(h))$. Indeed $|P(hx)|^2\ll h^4 x^4$ and $0\le 1-x^4/(1/4+x^2)^2\ll(1+x^2)^{-1}$. We omit atoms at zero in the arithmetic applications.

\section{Shared zeros and the identity residue class}
Fix $q\ge1$, write $G=(\mathbb Z/q\mathbb Z)^\times$, $n=\varphi(q)$, and let $\chi$ run over the $n$ characters modulo $q$. At a signed ordinate $\gamma$ write $m_\chi(\gamma)$ for the multiplicity of the primitive inducing $L$-function, with value zero if absent. The omitted Euler factors have no zeros in the open critical strip. Each sum over $\gamma$ uses distinct ordinates. Put $M(\gamma)=\sum_\chi m_\chi(\gamma)$.

For $u\in G$ define the residue-class prime-power signals by
\begin{align*}
 S_u(x)&=\sum_{\substack{k\le x\\k\equiv u\bmod q}}\Lambda(k)(1/k-1/x),&
 c_u&=\lim_{x\to\infty}(S_u(x)-n^{-1}\log x),\\
 E_u(t)&=e^{t/2}\big(S_u(e^t)-t/n-c_u\big),&
 F_\chi(t)&=\sum_{u\in G}\chi(u)E_u(t).
\end{align*}
For $q=1$, $E_1=E$. Under Dirichlet GRH the article's ramp explicit formula gives
\begin{equation}\label{eq:ef}
 F_\chi(t)=\sum_\gamma\frac{m_\chi(\gamma)e^{i\gamma t}}{1/4+\gamma^2}
       +O_q((1+t)e^{-t/2})\qquad(t\to\infty).
\end{equation}
Its Mellin integrand before normalization is $(-L'/L)(s,\chi)x^{s-1}/[s(s-1)]$. The residue at a strip zero is $m_\chi(\gamma)x^{-1/2+i\gamma}/(1/4+\gamma^2)$, and the pole at $1$ gives the logarithm and constant for each character, with the logarithm present only for the principal character. Shifting the contour to $\Re s=-1/2$ for the primitive inducer leaves a possible pole at $0$ contributing $O_q((1+\log x)/x)$. The omitted Euler factors contribute constants and errors of the same order by summing geometric tails. The local zero count $\sum_{y\le\gamma<y+1}m_\chi(\gamma)\ll_q\log(2+|y|)$ gives absolute convergence of the series. It also bounds the pair-product mass in a unit bin by $O_q(\log^2(2+|y|))$.

Let $\langle f,g\rangle=\lim_{Y\to\infty}Y^{-1}\int_0^Y f(t)\ol{g(t)}\dd t$ and put
\[
 W_{\chi\psi}(h)=\langle D_hF_\chi,D_hF_\psi\rangle,\qquad
 C_{uv}(h)=\langle D_hE_u,D_hE_v\rangle .
\]
The means of the differences vanish. The decaying error in \eqref{eq:ef} contributes zero to each mean, and unequal frequencies have zero mean product. Absolute convergence permits termwise averaging, so
\begin{equation}\label{eq:W}
 W_{\chi\psi}(h)=\sum_{\gamma\ne0}
 \frac{|P(h\gamma)|^2m_\chi(\gamma)m_\psi(\gamma)}{(1/4+\gamma^2)^2}\ge0 .
\end{equation}

\begin{theorem}[Pairwise inversion]\label{th:pair}
Under Dirichlet GRH, for any two characters modulo fixed $q$ and any finite $d\ge0$,
\[
 \frac{W_{\chi\psi}(h)}{\tau(h)}\longrightarrow B_{a,b}d
 \quad\Longleftrightarrow\quad
 \frac{\pi}{T\log T}\sum_{0<|\gamma|\le T}m_\chi(\gamma)m_\psi(\gamma)\longrightarrow d .
\]
The left limit may be taken on $h=1/k$. In particular $W_{\chi\psi}=o(\tau)$ precisely when the weighted number of common ordinates is $o(T\log T)$. This count folds the positive and negative ordinates, including for nonreal characters.
\end{theorem}
\begin{proof}
Apply Theorem~\ref{th:inverse} with mass $m_\chi(\gamma)m_\psi(\gamma)$ at $|\gamma|$ for each $\gamma\ne0$. The unit-bin bound verifies the denominator replacement. Taking $q=1$ gives the arbitrary-$\kappa$ part of Theorem~\ref{th:zeta} by conjugate symmetry.
\end{proof}

For the full matrix write $H_{u\chi}=\ol{\chi(u)}/n$. Finite Fourier inversion gives $C=HWH^*$ and $HH^*=I/n$. Thus $C/\tau$ converges entrywise precisely when all the folded pair counts in Theorem~\ref{th:pair} converge. If $D$ is the matrix of their limits, $C(h)/\tau(h)\to B_{a,b}HDH^*$.

\begin{theorem}[The identity-class criterion]\label{th:identity}
Under Dirichlet GRH the following are equivalent at any fixed modulus.
\begin{enumerate}[label=\textup{(\roman*)}]
\item $C_{11}(h)/\tau(h)\to B_{a,b}/n$ for one fixed $a,b>0$.
\item $\displaystyle\sum_{0<|\gamma|\le T}M(\gamma)(M(\gamma)-1)=o(T\log T)$.
\item For every $a,b>0$, $C(h)/\tau(h)\to (B_{a,b}/n)I$.
\end{enumerate}
In (i) it suffices to use $h=1/k$. Furthermore $\liminf C_{11}/\tau\ge B_{a,b}/n$.
\end{theorem}
\begin{proof}
Let $J_\chi(h)$ denote the right side of \eqref{eq:W} with the product replaced by $m_\chi(\gamma)$. Riemann--von Mangoldt for the fixed primitive inducer gives $\sum_{0<|\gamma|\le T}m_\chi(\gamma)\sim T\log T/\pi$, so $J_\chi/\tau\to B_{a,b}$. Put $C^0=H\diag(J_\chi)H^*$; then $C^0/\tau\to B_{a,b}I/n$. All entries of $W-\diag(J_\chi)$ are nonnegative. The diagonal entries count $m_\chi(m_\chi-1)$ with nonnegative weights, while the others count $m_\chi m_\psi$. Every entry of $H$ has modulus $1/n$, and its identity row is constant. This gives the finite-lag inequality
\begin{align}
 |C_{uv}-C^0_{uv}|&\le C_{11}-C^0_{11}\notag\\
 &=\frac{1}{n^2}\sum_{\gamma\ne0}\frac{|P(h\gamma)|^2M(\gamma)(M(\gamma)-1)}{(1/4+\gamma^2)^2}.\label{eq:dom}
\end{align}
This proves the lower bound. Theorem~\ref{th:inverse} for $d=0$, applied to the positive excess measure in \eqref{eq:dom}, gives (i)$\Leftrightarrow$(ii). The unit-bin mass is $O_q(\log^2(2+|y|))$ as before. Applying \eqref{eq:dom} for each geometry gives (ii)$\Rightarrow$(iii); taking the identity entry gives (iii)$\Rightarrow$(i).
\end{proof}

For $q=4$ the two primitive inducers are zeta and the Dirichlet beta function. At a given ordinate let their multiplicities be $m$ and $r$. Their combined excess is
\[
 (m+r)(m+r-1)=m(m-1)+r(r-1)+2mr .
\]
Thus the ordinary second-difference coefficient $1/3$ in the single class $1\bmod4$ is equivalent to all three weighted excesses being $o(T\log T)$. It forces coefficient $1/3$ for the other class $3\bmod4$ and zero limiting cross coefficient.

\section{The lower endpoint detects density-one simplicity}
For a fixed nonzero real contrast with distinct nodes $u_i$ and
$\sum_i c_i=\sum_i c_i u_i=0$, set
\[
 p(x)=\sum_i c_i e^{iu_i x},\qquad
 B=\frac1\pi\int_0^\infty\frac{|p(x)|^2}{x^4}\,dx.
\]
Write $V_c(h)$ for its long-time mean square on the shifted prime error. Under RH the following statements are equivalent.
\begin{enumerate}[label=\textup{(\roman*)}]
\item Simple zeros account for density one of all zeta zeros, counted with multiplicity.
\item $M_2(T)/N(T)\longrightarrow1$.
\item $V_c(h)/(h^3\log(1/h))\longrightarrow B$.
\end{enumerate}
This equivalence holds for \emph{every individual} nonzero contrast of the stated kind. For $(u_i)=(0,1,2)$, $(c_i)=(1,-2,1)$, the limiting constant is $2/3$. The proposition following the article's coalescence theorem gives the forward variance formula under a multiplicity asymptotic. Its remark that density-one simplicity would be insufficient overlooks a known uniform bound for the multiplicities of actual zeta zeros.

\paragraph{Multiplicity tails.}
Write $D_j(T)$ for the number of distinct zeros of multiplicity $j$ up to height $T$. Korolev's Theorem~A \cite{Korolev} gives absolute $A,c>0$ such that $D_j(T)\le A N(T)e^{-cj}$ for every $j\ge1$ and all sufficiently large $T$, uniformly in $j$. His convention at zero ordinates gives the same bound for right-continuous counts by taking a right limit. Set
\[
 d_T=1-D_1(T)/N(T),\qquad e_T=M_2(T)/N(T)-1 .
\]
For every integer $K\ge2$,
\[
 d_T\le e_T\le Kd_T+A\sum_{j>K}j^2 e^{-cj}.
\]
Indeed, the excess counts weight $j(j-1)$, whereas $N(T)-D_1(T)$ weights $j$ for $j\ge2$. Separating at $K$ gives the inequalities. Taking $T\to\infty$ and then $K\to\infty$ proves the equivalence of the first two conditions. Choosing $K=\lceil(2/c)\log(1/d_T)\rceil$ for $d_T$ sufficiently small yields $e_T=O(d_T\log(e/d_T))$.

\paragraph{Positivity at the lower endpoint.}
Use the spectral formula for $V_c(h)$, replacing $m_\gamma^2$ by $m_\gamma$ to define a baseline $V_c^{(1)}(h)$. The Riemann--von Mangoldt formula, with precisely the same change of scale and tail truncation as for the article's direct implication, gives $V_c^{(1)}(h)/(h^3\log(1/h))\to B$. The difference $V_c(h)-V_c^{(1)}(h)$ has nonnegative spectral weights $m_\gamma(m_\gamma-1)$. If (ii) holds, the normalized excess measures converge vaguely to zero, and the same tail control gives (iii).

Conversely, (iii) implies that the normalized excess above the baseline tends to zero. The exponential polynomial $p$ is analytic and nonzero by linear independence of exponentials with distinct frequencies. Choose a fixed $z>0$ such that $p$ has no zero on $[z/2,z]$; its zeros cannot accumulate at the origin. On $T/2<\gamma\le T$ put $h=z/T$. Then $|p(h\gamma)|^2$ is bounded below by a positive constant, and $(1/4+\gamma^2)^{-2}$ is bounded below by a constant times $T^{-4}$. Hence the vanishing normalized excess gives
\[
 \sum_{T/2<\gamma\le T}m_\gamma(m_\gamma-1)=o(T\log T).
\]
Summing dyadically proves $M_2(T)-N(T)=o(T\log T)$. For explicit control, fix $\eta>0$ and bound every dyadic block with upper endpoint $Y\ge Y_0(\eta)$ by $\eta Y\log Y$. The sum is at most $2\eta T\log T$ plus a bounded initial segment. This proves (ii) by the zero-counting asymptotic. Positivity makes this endpoint argument valid even for a contrast whose Mellin transform has a real zero away from the origin.

\paragraph{A dyadic observation already suffices.}
In condition (iii) it suffices to take $h=2^{-n}$. More generally, fix any strictly decreasing sequence $h_n\downarrow0$ with $h_n/h_{n+1}\le R<\infty$. Choose $z>0$ with $p(x)\ne0$ throughout $(0,z]$. For each large $T$, at the first index $n$ such that $h_n\le z/T$ we have $z/R<h_nT\le z$. Throughout $T/2<\gamma\le T$ the arguments $h_n\gamma$ belong to the compact interval $[z/(2R),z]$, on which $|p|$ has a positive minimum. Since $h_n^3\log(1/h_n)$ is comparable to $T^{-3}\log T$, the preceding shell estimate and its dyadic sum go through verbatim. Thus a dyadic sequence of long-time variances of a single ordinary second difference contains the full density-one simplicity criterion.
\section{A smaller upper bound for every coalescing contrast}
\label{sec:local-pair}
Fix $0<a<b$. Under RH we have
\[
  b-a\ \le\ \liminf_{T\to\infty}
  \frac{M_2(bT)-M_2(aT)}{T\log T/(2\pi)}
  \ \le\ \limsup_{T\to\infty}
  \frac{M_2(bT)-M_2(aT)}{T\log T/(2\pi)}
  \ \le\ 1.3208(b-a).
\]
This local bound answers the article's question about the factor $4/3$ in its zeta coalescence theorem. Every limiting density of the measures
$2\pi(T\log T)^{-1}\sum_{\gamma>0}m_\gamma^2\delta_{\gamma/T}$
lies between $1$ and $1.3208$ almost everywhere. Thus each upper comparison with $\tfrac43\mathcal B$ for coalescing quadratic forms improves to $1.3208\mathcal B$. The same substitution holds for the small eigenvalues and the optimal linear prediction error by the article's block elimination. For the second difference, the upper bound becomes $(2/3)1.3208=0.8805333\ldots$, compared with $8/9$.

\paragraph{Restricting pair correlation to a macroscopic height interval.}
Use $w(y)=4/(4+y^2)$ and $D_T=T\log T/(2\pi)$. In the following double sums each zero is repeated according to its multiplicity. For $I=(aT,bT]$ set
\[
  F_{I,T}(v)=D_T^{-1}\sum_{\gamma,\gamma'\in I}
            e^{iv(\gamma-\gamma')\log T}w(\gamma-\gamma').
\]
This function is nonnegative for every real $v$, since $w(y)=\int_{\mathbb R}e^{-2|s|}e^{isy}\,ds$ turns the double sum into an integral of an absolute square. On each compact subinterval of $(-1,1)$ the measures $F_{I,T}(v)\,dv$ converge weakly against compactly supported tests to
\[
 (b-a)(\delta_0+|v|\,dv).
\]
Indeed, first take the prefix $(0,cT]$ for $c=a,b$. Montgomery's pair-correlation theorem \cite{mont}, in its usual normalization at height $X=cT$, gives
$(1+o(1))\log X\,X^{-2|u|}+|u|+o(1)$ for the normalized form factor, uniformly on $|u|\le1-\delta$. Change variables to $u=v\log T/\log X$ and multiply by $D_X/D_T\to c$. The peak $\log X\exp(-2|v|\log T)$ has mass tending to one. Subtracting the two prefix formulas leaves, in addition to the desired double sum, pairs lying on opposite sides of $aT$. Their unnormalized total absolute contribution is $O_{a,b}(\log^3 T)$ uniformly in $v$. To see this, divide each side of $aT$ into distance bins $[j,j+1)$ and $[k,k+1)$. Each bin contains $O_b(\log T)$ zeros, and the sum of $(1+(j+k)^2)^{-1}$ over $0\le j,k\le bT+1$ is $O_b(\log T)$. Thus the cross pairs disappear after normalization.

\paragraph{The test-function inequality.}
Write $\widehat f(v)=\int f(x)e^{-2\pi ivx}\,dx$. Take a real even continuous $f\ge0$ with $f(0)=1$, with $f,\widehat f\in L^1$, and with $\widehat f(v)\le0$ for $|v|\ge1$. The diagonal pairs of equal ordinate contribute $M_2(bT)-M_2(aT)$. For $0<d<1$ put $f_d(x)=f(dx)$. By positivity in the ordinate variable and Fourier inversion,
\[
 \frac{M_2(bT)-M_2(aT)}{D_T}
 \le \int_{\mathbb R}\widehat f_d(v)F_{I,T}(v)\,dv
 \le \int_{-d}^d\widehat f_d(v)F_{I,T}(v)\,dv.
\]
The second inequality uses both $F_{I,T}\ge0$ and the sign of $\widehat f_d$ outside $[-d,d]$. The local weak convergence applies also to this sharp cutoff, whose endpoints have zero mass in the limiting measure. Taking $T\to\infty$, then $d\uparrow1$, gives the upper bound $(b-a)Q(f)$, where
\[
 Q(f)=\widehat f(0)+\int_{-1}^1 |v|\widehat f(v)\,dv .
\]
Chirre, Gon\c calves and de Laat construct a test with $Q(f)\le1.3208$ in the proof of their Theorem~1 \cite{ChirreLP}. In their verification convention, write $h$ for their even polynomial times a Gaussian, with $h(x)\le0$ for $|x|\ge R$ and $\widehat h\ge0$. Our test is $f(x)=\widehat h(x/R)/\widehat h(0)$, so $Q(f)$ equals their rigorously bounded quantity $Z(h)$ in Section~4.1.1. The lower bound follows from $m^2\ge m$ and the Riemann--von Mangoldt formula. Applying these interval inequalities to continuous compactly supported nonnegative tests proves the claimed bounds on all vague limits. Truncation for the spline kernels follows from the article's bounds at zero and infinity. This argument transfers a published test to the local covariance problem. A bound for $M_2(T)$ alone would not bound these differences by subtraction.

\paragraph{A closed-form certificate.}
An independently evaluable choice gives the classical Montgomery--Taylor constant
\[
 Q_{\rm MT}=\tfrac12+\tfrac1{\sqrt2}\cot(1/\sqrt2)=1.3274992963\ldots .
\]
On $J=[-1/2,1/2]$, choose $g(t)=\cos(\sqrt2 t)/(\sqrt2\sin(1/\sqrt2))$, zero elsewhere, and $f=|\widehat g|^2$. Then $\int g=1$, $f(0)=1$, and $\widehat f$ is supported in $[-1,1]$. Direct Fourier calculation yields
\[
 Q(f)=\int_Jg(t)\left[g(t)+\int_J |t-u|g(u)\,du\right]dt.
\]
The bracket is constant on $J$ because its second derivative equals $g''+2g=0$ in the interior and it is even. Its value at $1/2$ is $g(1/2)+1/2=Q_{\rm MT}$. An alternative using only rational integrals takes $g(t)=1+z(1/12-t^2)$ on $J$; here $Q(f)=4/3-z/90+z^2/189$, giving $531/400$ at $z=21/20$. Already this elementary certificate proves strict improvement. The optimized value over all possible arithmetic methods remains open.
\section{The cubic rate in every residue-class direction}\label{sec:familyrate}

Fix distinct primitive Dirichlet characters $\chi_1,\ldots,\chi_d$,
allowing the conductor-one character. At each distinct real ordinate
$\gamma$ write $m(\gamma)=(m_1(\gamma),\ldots,m_d(\gamma))^{\mathsf T}$,
where $m_j$ is the multiplicity of $1/2+i\gamma$ for $L(s,\chi_j)$.
The following estimate controls cancellations between different zero divisors.
It uses Dirichlet GRH. Throughout, $A,B$ are fixed with $0<A<B$.

\begin{theorem}[Local bounds for the multiplicity matrix]\label{thm:familyrate}
There are constants $0<c<C<\infty$ depending only on the family such that,
for all sufficiently large $T$ (depending also on $A,B$),
\[
 c(B-A)T\log T\, I_d \ \preceq\!
 \sum_{AT<\gamma\le BT}m(\gamma)m(\gamma)^{\mathsf T}
 \ \preceq C(B-A)T\log T\, I_d .
\]
Here $\preceq$ denotes the ordering of Hermitian quadratic forms. In particular the
estimate is uniform in complex coefficient vectors. It holds on negative-height
intervals by conjugating the family.
\end{theorem}

\begin{proof}
Put $L=\log T$ and choose a nonzero smooth function $w$ compactly supported
in $(a,b)$, with $0<a<b<1/8$. Set
$f(x)=\int_{\mathbb R}w(v)e^{ivx}\,dv$. For a unit vector $c\in\mathbb C^d$
write $z(\gamma)=\sum_j c_jm_j(\gamma)$, $U(n)=\sum_jc_j\chi_j(n)$ and
$X(t)=\sum_\gamma z(\gamma)f(L(\gamma-t))$.
The Dirichlet explicit formula
(\cite[Lemma 2.1]{liuye}, applied to each primitive inducer) gives, uniformly
on compact positive ranges of $t/T$,
\[
 X(t)=-\frac1L\sum_{n\le T^b}\frac{\Lambda(n)U(n)}{\sqrt n}\,
             w(\log n/L)n^{-it}+o(1).
\]
Indeed the Fourier test is $L^{-1}w(v/L)e^{-itv}$, supported on positive
frequencies, so only one prime sum occurs. The archimedean density is
$O_j(\log(2+|t|))$ with derivative $O_j(1/|t|)$ at large height.
After scaling the archimedean integral, its constant term vanishes by
$\int f=2\pi w(0)=0$; Stirling's formula and the Schwartz tails bound
the difference by $o(1)$. The pole terms for the conductor-one character
are $f(L(\pm i/2-t))=o(1)$ by repeated integration by parts in $v$:
their possible factor $T^{b/2}$ is dominated by a sufficiently high
power of $t$. All errors are uniform in $c$ by finite dimensionality.

Use normalized integration on any interval $J$ of length $H\asymp_{A,B}T$
within this range. For $D(t)=\sum_{n\le N}d_n n^{-it}$, direct integration yields
\[
 H^{-1}\int_J|D(t)|^2dt=
 \bigl(1+O(N\log(2N)/H)\bigr)\sum_n|d_n|^2 .
\]
Indeed $1/|\log(m/n)|\le N/|m-n|$ and
$2|d_m d_n|\le |d_m|^2+|d_n|^2$.
For the prime part $P$ of the displayed polynomial including $1/L$,
the prime number theorem on reduced classes modulo the lcm of the
conductors, partial summation and character orthogonality give
\[
 \sum_p |d_p|^2\longrightarrow
       \int_a^b v|w(v)|^2dv>0 .
\]
The distinct primitive characters induce distinct characters at this common
modulus. Entrywise convergence of the finite Gram matrix makes the limit
uniform on the unit sphere. Higher prime powers have coefficient square sum
$O_d(L^{-2})$ since $\sum_{p,k\ge2}(\log p)^2/p^k<\infty$.
Consequently $\|X+P\|_{L^2(J,dt/H)}=o(1)$.
Apply the same mean-value estimate to $P^2$, of length at most $T^{2b}$.
Unique factorization gives
$\sum_n|\sum_{p p'=n} d_pd_{p'}|^2\le2(\sum_p|d_p|^2)^2$.
Thus $\|P\|_4=O(1)$ and $\|P\|_2$ is bounded below.
Interpolation $\|P\|_2\le\|P\|_1^{1/3}\|P\|_4^{2/3}$ proves
$H^{-1}\int_J|X|\,dt\gg 1$, with a lower constant independent of $A,B$.

Take $J$ to be the middle half of $[AT,BT]$. The fixed-conductor count,
with multiplicity, in each unit interval at height $y$ is
$O_d(\log(2+|y|))$. Thus the contribution to $X(t)$ from outside
$[AT,BT]$ is uniformly $o(1)$ on $J$ by rapid decay; an endpoint
can likewise be omitted. Every remaining ordinate contributes at most
$|z(\gamma)|\|f\|_1/L$ to the integral. Hence
$\sum_{AT<\gamma\le BT}|z(\gamma)|\gg (B-A)TL$.
The Riemann--von Mangoldt count gives at most $O_d((B-A)TL)$ ordinates
in that interval, with an implied constant independent of its fixed
endpoints for sufficiently large $T$. Cauchy--Schwarz yields the lower bound.

For the upper bound choose a nonnegative Schwartz test $g$ with smooth
Fourier transform supported in $(-b,b)$ and $g(0)>0$, for instance the
square modulus on the real axis of an inverse Fourier bump of width $b/2$
and nonzero integral. The same explicit formula for
$Y_j(t)=\sum_\gamma m_j(\gamma)g(L(\gamma-t))$ now gives two
short polynomials and an $O(1)$ archimedean term. Each polynomial has
bounded coefficient square sum by
$L^{-2}\sum_{p\le T^b}(\log p)^2/p=O(1)$ and the preceding bound for
higher powers. The pole term is still $o(1)$. Hence
$\int_{AT-1}^{BT+1}|Y_j|^2dt\ll (B-A)T$.
Since all summands are nonnegative, the diagonal contribution bounds
this integral below by
$L^{-1}\int_{-1}^1g(u)^2du\ \sum_{AT<\gamma\le BT}m_j(\gamma)^2$.
Summing over $j$ and using
$|z(\gamma)|^2\le\sum_jm_j(\gamma)^2$ gives the upper bound.
\end{proof}

\paragraph{Application to the ramps.}
Take the $d=\varphi(q)$ primitive inducers of the characters modulo $q$.
Denote by $Z_a(t)$ the centered normalized ramp for the reduced class
$a$, that is, $Z_a=E_a-\langle E_a,1\rangle$, as in the source ramp formula (label
\texttt{zer:lem:ramp}). For distinct fixed offsets $u_1,\ldots,u_r$,
let $b_j\in\mathbb R^d$ be vectors of residue-class coefficients satisfying
$\sum_jb_j=\sum_j u_jb_j=0$. Define
$H_b(v)=\sum_j b_j(v-u_j)_+$. Writing $\mathcal M$ for long-time mean, one has
\[
 \mathcal M\left|\sum_j b_j^{\mathsf T} Z(t+h u_j)\right|^2
 \ \asymp_{q}\ h^3\log(1/h)\int_{\mathbb R}\|H_b(v)\|^2dv
 \qquad (h\downarrow0).
\]
The comparison holds uniformly in $b$ for fixed offsets. In particular
the $d$ eigenvalues of the residue-class second-difference covariance all
have order $h^3\log(1/h)$, even with coincidences among ordinates.

To verify the comparison, scale both signs of the matrix measure in
Theorem~\ref{thm:familyrate} by placing $m(\gamma)m(\gamma)^{\mathsf T}/(T\log T)$
at $\gamma/T$, where $T=1/h$. The theorem bounds every vague limit
between $c I_d\,dx$ and $C I_d\,dx$ on each open half-line: test
constant vectors on intervals, take limits using continuous step
approximations, then polarize. The upper bound for $[T,2T]$ summed
over dyadic intervals gives $\sum_{0<|\gamma|\le X}\|m(\gamma)\|^2
\ll_q X\log(2+X)$ for $X\ge1$. This controls the ends of the integral.
Indeed for $p(x)=\sum_j b_j e^{iu_jx}$,
$\|p(x)\|\ll_u (\sum_j\|b_j\|^2)^{1/2}\min(x^2,1)$.
The mass near zero has limiting upper bound $O_q(\delta)$ on
$[-\delta,\delta]$; the weighted tail with weight $|x|^{-4}$
on $|x|>R$ is $O_q((1+\log R)R^{-3})$ for $T\ge2$.
The character-to-class matrix
$D_{a\chi}=\overline{\chi(a)}/d$ satisfies $DD^*=I_d/d$.
The ramp expansion and Parseval therefore bound the limiting normalized
variance above and below by constant multiples of
$d^{-1}\int_{\mathbb R}\|p(x)\|^2x^{-4}dx$:
the exact denominator after scaling is $(x^2+1/(4T^2))^2$.
The latter integral equals $2\pi\int\|H_b(v)\|^2dv$ by Plancherel.
This form is positive definite on the contrast space, since $H_b''$
has the coefficients $b_j$ as its atoms. Compactness of its unit sphere
gives uniformity. For offsets $(0,1,2)$ use $(b_1,b_2,b_3)=(b,-2b,b)$.

For two functions and the difference of their multiplicities, Fujii and
Bombieri--Perelli \cite{bp} already obtained lower bounds of order $T\log T$.
The role of the matrix estimate here is its uniformity across all
direction vectors of a fixed family and its transfer to every
affine-annihilating residue-class contrast. The argument gives the scale
of the covariance; its exact matrix of leading constants contains
additional information about coincident zeros.

\paragraph{All eigenvalue orders.}
For $r\ge3$ the $rd$-dimensional joint covariance of
$(Z(t+h u_1),\ldots,Z(t+h u_r))$ has $d$ eigenvalues of order $1$,
$d$ of order $h^2$ and $(r-2)d$ of order $\tau=h^3\log(1/h)$.
In particular its determinant has order $h^{(3r-4)d}\log(1/h)^{(r-2)d}$.

Indeed use an orthonormal offset basis beginning with the constant vector
and the centered vector $(u_j-\bar u)_j$ normalized to length one.
The remaining basis vectors annihilate affine functions.
Realize covariance as the Gram matrix of Fourier coefficient vectors in
$\ell^2$ of the distinct noncentral ordinates, using both signs.
For each residue class, the first observation vector converges in $\ell^2$
to a nonzero constant multiple of its ramp vector. After division by $h$,
the second converges to a nonzero constant multiple of its derivative
vector, by $|e^{iy}-1|\le |y|$ and convergence of
$\sum_{\gamma\ne0}|\gamma|^2\|m(\gamma)\|^2/(1/4+\gamma^2)^2$.
Their joint limiting Gram is positive definite. In fact any relation would
give $m(\gamma)^{\mathsf T}(v+i\gamma w)=0$ for all noncentral ordinates
and some fixed complex $v,w$ after the invertible character transform.
On $[T,2T]$ this would imply
$T^2\sum|m(\gamma)^{\mathsf T}w|^2\le\sum|m(\gamma)^{\mathsf T}v|^2$;
Theorem~\ref{thm:familyrate} forces $w=0$ and then $v=0$.
The remaining observation vectors divided by $\sqrt\tau$ have Gram matrices
uniformly bounded above and below by the contrast estimate. Each such
vector converges weakly to zero: it is bounded in $\ell^2$, and its
coefficient at any fixed ordinate is $O_\gamma(h^2/\sqrt\tau)=o(1)$.
The mixed Gram entries against the strongly convergent vectors therefore
vanish. The normalized full Gram lies between positive constant multiples
of the identity for small $h$. Undoing the diagonal normalization gives
the asserted eigenvalue orders by the min--max principle.

The multiplicity estimate itself extends to all intervals $[T,T+H]$ with
$T^\delta\le H\le T$ for each fixed $0<\delta\le1$, with constants
depending also on $\delta$. Choose $2b<\delta$ in the proof.
Then $T^{2b}\log T/H=o(1)$ uniformly, the middle half stays at distance
$\gg H$ from the endpoints, and the zero-count and Schwartz-tail estimates
used there are uniform. The bounds in that case are multiples of $H\log T$.

\paragraph{Ordinary simple zeros of the cyclotomic field.}
For the family of all primitive inducers modulo $q$, the three equivalent
conditions in Theorem~\ref{th:identity} are equivalent to one more:
simple zeros have density one among the nontrivial zeros of the Dedekind
zeta function of $\mathbb Q(\zeta_q)$, counted with multiplicity. It suffices
in the identity-class condition to sample $h=2^{-k}$.

Here are the details about large multiplicities. On $[T-1,2T+1]$ the
nonnegative test count $Y_j$ used in Theorem~\ref{thm:familyrate} has
bounded normalized fourth moment. Each prime polynomial in it has bounded
fourth moment by the estimate for $P^2$ above. The higher powers are
bounded pointwise, since
\[
 \frac1L\sum_{p^k\le T^b,\ k\ge2}\frac{\log p}{p^{k/2}}
 \ll \frac1L\left(\sum_{p\le T^b}\frac{\log p}{p}
       +\sum_p\frac{\log p}{p^{3/2}(1-p^{-1/2})}\right)\ll1.
\]
The pole and archimedean terms are bounded as before. Positivity in
$Y_j^4$ gives $\sum_{T<\gamma\le2T}m_j(\gamma)^4\ll_q T\log T$,
using $L^{-1}\int_{-1}^1g(u)^4du$ from each diagonal term. Sum over
dyadic intervals and both signs. Since $(\sum_{j=1}^dm_j)^4\le
d^3\sum_jm_j^4$, this gives
$\sum_{0<|\gamma|\le T} M(\gamma)^4\ll_q T\log T$.
The cyclotomic factorization into the primitive Dirichlet $L$-functions
makes $M(\gamma)$ the multiplicity of that Dedekind zeta zero.
Its total count here is asymptotic to $dT\log T/\pi$.
If its simple zeros have density one, the number of distinct ordinates
of multiplicity at least two is $o(T\log T)$; Cauchy--Schwarz with the
fourth-moment bound proves that their total squared multiplicity is
$o(T\log T)$. Conversely $M\le M(M-1)$ on $M\ge2$, so the
excess condition gives density-one simplicity.

Finally suppose the identity-class variance attains the minimum along
$h=2^{-k}$ alone. Subtract its linear-multiplicity baseline in
\eqref{eq:dom}. For dyadic $T=2^j$ select $h=2^{-j-s}$, where the
positive integer $s$ is fixed sufficiently large that $|P|$ is bounded
away from zero on $[2^{-s-1},2^{-s}]$ (it has an isolated zero at the
origin). The nonnegative excess from $T/2<|\gamma|\le T$ is at least
a positive constant times $T^{-4}\sum M(\gamma)(M(\gamma)-1)$ on
that shell. Vanishing normalized excess along the sampled sequence
therefore bounds this sum by $o(T\log T)$. Dyadic summation, then
monotonicity between dyadic endpoints, give condition~(ii) in
Theorem~\ref{th:identity}. The reverse follows from the full limit.

\begin{corollary}[A criterion that includes the critical-line assertion]
Fix $q$ and let $E_1$ be the identity-class ramp. For $a=b=1$ the following two assertions are equivalent, without assuming GRH beforehand.
\begin{enumerate}[label=\textup{(\roman*)}]
\item The mean squares exist for all sufficiently large $k$ and satisfy \[ C_{11}(2^{-k})/\tau(2^{-k})\longrightarrow 2/(3\varphi(q)). \]
\item Every nontrivial zero of $\zeta_{\mathbb Q(\zeta_q)}$ lies on the critical line and its simple zeros have density one, counted with multiplicity.
\end{enumerate}
\end{corollary}
\begin{proof}
Use the unconditional centering constant
$c_1=n^{-1}\sum_{\chi\bmod q}(\operatorname{FP}_{s=1}(-L'/L)(s,\chi)-\delta_\chi)$,
where $\operatorname{FP}$ takes the constant Laurent coefficient and
$\delta_\chi$ indicates the principal character. This agrees with the prime-sum limit defining $c_1$ by the fixed-modulus prime number theorem with its usual integrable remainder, or by partial summation in each residue class followed by Abelian passage to $s=1$.
For $\Re s>1/2$ set $\lambda=s-1/2$. Termwise Laplace transformation gives
\[
 {\cal L}E_1(s)=
 \frac{\sum_{\chi\bmod q}(-L'/L)(1+\lambda,\chi)}
 {n\lambda(1+\lambda)}-\frac{1}{n\lambda^2}-\frac{c_1}{\lambda}.
\]
In fact the transform of each ramp summand before centering contains the integral from $\log k$ to infinity of $(1/k-e^{-t})e^{-\lambda t}$, equal to $k^{-1-\lambda}/[\lambda(1+\lambda)]$.
For the second difference the transform is $(e^{hs}-1)^2{\cal L}E_1(s)$ plus an entire function, by shifting the lower integration endpoint.
A finite upper long-time mean square at any one $h>0$ makes this last transform holomorphic on $\Re s>0$: the integral of $|D_h E_1|^2$ up to $Y$ is $O(1+Y)$, and Cauchy--Schwarz gives exponential integrability locally uniformly on that half-plane.
At a strip zero $\rho$ with $\Re\rho>1/2$, the sum of logarithmic derivatives has a pole at $s=\rho-1/2$ with nonzero residue, since multiplicities add. Neither the centering terms nor the difference multiplier cancels it. Meromorphic continuation and the identity theorem therefore exclude such zeros. The functional equation and cyclotomic factorization give the critical-line assertion in (ii). The density-one part and the reverse implication now follow from the preceding dyadic identity-class theorem, whose proof uses GRH only for the primitive inducers of $q$.
\end{proof}

\section{A contrast that loses information}
The converse for arbitrary limiting coefficients uses the strict inequality $(a^2+b^2)/(a+b)^2<1$ for three-point differences. A more complicated contrast can hide fluctuations in a general positive spectral measure. On $x\ge1$ consider
\[
 \dd A(x)=\log(2+x)(1+\epsilon\cos(\omega\log x))\dd x,
 \qquad 0<\epsilon<1,\quad \omega=\pi/\log2 .
\]
Substitution $x=Ty$ and integration of the two complex powers give
\[
 A(T)/(T\log T)=1+\epsilon\Re\{T^{i\omega}/(1+i\omega)\}+o(1).
\]
This has no limit. Take $a=b=1$. Substitution $y=hx$ shows that $U(h)/\tau(h)$ equals $\int K(y)\dd y+\epsilon\Re\{h^{-i\omega}\int K(y)y^{i\omega}\dd y\}+o(1)$. Indeed $K(y)(1+|\log y|)$ is integrable, so dominated convergence controls the slowly varying logarithm. Replacing $h$ by $2h$ reverses the oscillation. Hence $(U(h)+U(2h)/8)/\tau(h)\to 2\pi B_{1,1}$.

This sum is the variance for one real five-point contrast. On $|z|=1$ there is a real quadratic $Q$ with $|Q(z)|^2=1+|1+z|^4/8$. Explicitly, put $s=(1+\sqrt3)/2$ and take the coefficients of $Q$, in ascending order, to be $(s+\sqrt s)/2,\ 1/(2s),\ (s-\sqrt s)/2$. Comparison of the three Laurent coefficients proves the identity. The coefficients of $(1-z)^2Q(z)$ annihilate constants and linear functions, and its squared modulus is $|1-z|^4+|1-z^2|^4/8$. Thus this five-point contrast can have a limiting coefficient while its normalized spectral count oscillates. The example concerns an abstract positive measure; its realization by zeta zeros is a separate question.

\paragraph{Scope and further questions.}
The source article's cubic coalescence section supplies the forward comparison by splines and the forward implication for a prescribed multiplicity asymptotic. The reverse implication here uses the three-point geometry; the identity-class theorem also uses positivity before character inversion. Inverse regular-variation problems and logarithmically periodic obstructions have a general theory \cite{inverse}. Cram\'er's mean-square theorem \cite{cramer} already relates a prime error to a convergent sum over zeros. Our limit probes the high zeros through a shrinking log-scale lag. It is insensitive to any fixed finite set of zeros. The distinct short-interval variance problems associated with full pair correlation retain contributions from pairs of different ordinates. Here the infinite-time average eliminates all unequal frequencies before the lag tends to zero.

The vendor manifest identifies upstream revision \href{https://github.com/lsernau/grh-lean/tree/72fead102f663d5c7a407964a13dc0cf9d977207}{\texttt{72fead102f66}} of the readable edition, a separate Lean package. The vendored file \texttt{grh/src/Main.lean} exports \texttt{riemannHypothesis} and \texttt{LFunction\_\allowbreak eq\_zero\_re}. This note uses the stated RH and Dirichlet GRH as inputs to a classical argument. The last corollary also proves the corresponding two-way implication with the critical-line statement included in the conclusion. The new proofs have not been formalized and we did not rebuild the upstream Lean development. Novelty here means a proved extension of the identified sources; historical priority remains to be established. Can a quantitative error term for the variance give an effective error term for the sharp cutoff $M_2(T)$? Which other individual stencils characterize arbitrary finite multiplicity limits?

\begin{thebibliography}{9}
\bibitem{repo} V. Reshetnikov, \emph{ProveIt}, revision \texttt{c0f2782360ddffd79a20c6d828d582e214a2b8c3}.
\href{https://github.com/VladimirReshetnikov/ProveIt/blob/c0f2782360ddffd79a20c6d828d582e214a2b8c3/Analysis/GeneralizedRiemannHypothesis/Article/sections/zer.tex}{Article, ramp formula and cubic coalescence section};
\href{https://github.com/VladimirReshetnikov/ProveIt/blob/c0f2782360ddffd79a20c6d828d582e214a2b8c3/Analysis/GeneralizedRiemannHypothesis/vendor/grh-lean/grh/src/Main.lean}{vendored Lean declarations}.
\bibitem{cramer} H. Cram\'er, \emph{Ein Mittelwertsatz in der Primzahltheorie}, Math. Z. 12 (1922), 147--153.
\bibitem{liuye} J. Liu and Y. Ye, \emph{Distribution of zeros of Dirichlet $L$-functions and the least prime in an arithmetic progression}, Acta Arith. 119 (2005), 13--38. \href{https://www.impan.pl/shop/publication/transaction/download/product/83731}{Lemma 2.1}.
\bibitem{bp} E. Bombieri and A. Perelli, \emph{Distinct zeros of $L$-functions}, Acta Arith. 83 (1998), 271--281.
\bibitem{mont} H. L. Montgomery, \emph{The pair correlation of zeros of the zeta function}, Proc. Sympos. Pure Math. 24 (1973), 181--193. \href{https://public.websites.umich.edu/~hlm/paircor1.pdf}{Author's copy}.
\bibitem{ChirreLP} A. Chirre, F. Gon\c calves and D. de Laat, \emph{Pair correlation estimates for the zeros of the zeta function via semidefinite programming}, Adv. Math. 361 (2020), 106926. \href{https://arxiv.org/abs/1810.08843}{arXiv:1810.08843}.
\bibitem{Korolev} M. A. Korolev, \emph{On multiple zeros of the Riemann zeta function}, Izv. Math. 70 (2006), 427--446. \href{https://doi.org/10.1070/IM2006v070n03ABEH002317}{Theorem A}.
\bibitem{inverse} M. Jacobsen, T. Mikosch, J. Rosi\'nski and G. Samorodnitsky, \emph{Inverse problems for regular variation of linear filters, a cancellation property for $\sigma$-finite measures and identification of stable laws}, Ann. Appl. Probab. 19 (2009), 210--242. \href{https://doi.org/10.1214/08-AAP540}{DOI}.
\end{thebibliography}
\end{document}
